\clearpage
\section*{四、实验心得与思考题}

\subsection*{1. 实验心得}

\begin{enumerate}
\item 拉脱法真正需要的是最大拉力与液膜破裂后拉力之差。用示波器记录整个拉脱过程、暂停后再用光标读峰值，比盯着数字表读数更可靠。
\item 定标数据的线性很好，增重和减重读数最多只差 1 mV，说明传感器本身不是主要误差来源。结果偏低主要来自水和圆环的洁净程度以及调零方式。
\item 虽然测量值整体偏低，但 $\delta$ 随温度升高而减小的规律和变化率都与标准值相符，温度对表面张力的影响得到了验证。
\item 如果重做实验，我会用镊子取放圆环，彻底清洗后在干燥状态下调零，并在每次拉脱后记录基线读数，以减小系统误差。

\end{enumerate}
\subsection*{2. 思考题}

\subsubsection*{1. 在圆环上提水膜即将破裂时 $F=mg+F_{\text{拉}}$ 成立，若过早读数，对实验结果会有什么影响？}

\begin{enumerate}
\item 液膜被拉起的过程中，拉力逐渐增大。只有当液膜接近竖直、即将破裂时，表面张力的竖直分量才达到最大，$F=mg+F_{\text{拉}}$ 才成立。
\item 若过早读数，液膜与竖直方向仍有夹角，$\cos\theta<1$，读到的拉力小于最大值，$\Delta U$ 偏小，由 $\delta=K\Delta U/[\pi(d_1+d_2)]$ 得到的 $\delta$ 也偏小。
\item 本实验用示波器记录拉脱全过程，暂停后用光标读最大值，可以避免过早读数。但也不能取破裂瞬间的读数，因为最大值之后拉力已略有下降。

\end{enumerate}
\subsubsection*{2. 圆环或玻璃容器不洁净，$\delta$ 会有什么变化？}

\begin{enumerate}
\item 油脂、灰尘等杂质多为表面活性物质，进入水中后聚集在液面，降低水的表面张力，测得的 $\delta$ 偏小。
\item 圆环表面有油污时，水对圆环的润湿变差，接触角增大，液膜容易提前从环上脱开，测得的最大拉力偏小，$\delta$ 同样偏小。
\item 本次测得的 $\delta$ 比标准值低约 23\%，水和圆环不够洁净是主要原因之一。因此，实验前应清洗圆环和器皿，并避免用手直接接触圆环。
\end{enumerate}
